% !TeX root = ../../Book4.tex
% 纲要标题：### 14.2 谱序列的构造
% 纲要位置：计划纲要.md，第 2833 行。

\section{谱序列的构造}
\label{b4:sec:1402}

第零页只看伴随分次，后续微分则测量代表元还能向深层提升多少次。本节先从代表元直接构造各页，再证明每一新页由前一页取同调得到。

% 纲要内容：
%
% * E_0、E_1 及微分
% * 页间同调关系
%

\subsection{\texorpdfstring{\(E_0\)、\(E_1\) 与微分}{E0、E1 与微分}}
\label{b4:subsec:1402:01}

设 \((C,F)\) 为模的递减滤过复形。第零页为
\[
 E_0^{p,q}=F^pC^{p+q}/F^{p+1}C^{p+q},\qquad \mathrm{d}_0[x]=[\mathrm{d}x].
\]
微分保持滤过，保证它良定义。因此
\[
 E_1^{p,q}=H^{p+q}(F^pC/F^{p+1}C).
\]
若 \(x\) 代表此页的类，便有 \(\mathrm{d}x\in F^{p+1}\)。再取 \(\mathrm{d}x\) 在下一层的类，就得到 \(\mathrm{d}_1[x]\)。

\ProseParagraphBreak
例如，取零次和一次的两项复形 \(C=(\mathbb Z\xrightarrow{2}\mathbb Z)\)，并令 \(F^pC=C\)（\(p\le0\)）、\(F^1C=(0\to\mathbb Z)\)、\(F^pC=0\)（\(p\ge2\)）。乘 \(2\) 把源的滤过层零送入靶的滤过层一，所以伴随分次上的 \(\mathrm{d}_0\) 为零。第零页与第一页都只有两个非零项，但第一页微分恢复了原复形中跨过一层滤过的映射：
\[
 E_1^{0,0}=\mathbb Z\xrightarrow{\ \mathrm{d}_1=\times2\ }
 E_1^{1,0}=\mathbb Z.
\]
于是第二页只剩 \(E_2^{1,0}=\mathbb Z/2\)，以后再无可能的非零微分。它与 \(H^0(C)=0\)、\(H^1(C)=\mathbb Z/2\) 相符：先逐层看时得到的两个自由群，通过后续微分才合成为原复形的挠同调。

\ProseParagraphBreak
为同时写出后面各页，令 \(n=p+q\)，并对 \(r\ge0\) 定义
\[
 Z_r^{p,q}=F^pC^n\cap \mathrm{d}^{-1}(F^{p+r}C^{n+1}).
\]
特别地 \(Z_0^{p,q}=F^pC^n\)。当 \(r\ge1\) 时，置
\begin{equation}\label{b4:eq:spectral-page-formula}
 E_r^{p,q}=
 \frac{Z_r^{p,q}}
 {Z_{r-1}^{p+1,q-1}+\mathrm{d}Z_{r-1}^{p-r+1,q+r-2}}.
\end{equation}
分母第一项容许更改较深一层的代表元，但其微分也要足够深；第二项容许减去已经可见的余边。两项确实包含于分子：第一项的微分位于 \(F^{p+r}\)，第二项位于 \(F^p\) 且再次微分为零。取 \(r=1\)，公式正是上面第零页的上同调。

\subsection{页间同调关系}
\label{b4:subsec:1402:03}

\begin{theorem}\label{b4:thm:filtered-spectral-sequence}
设 \(R\) 为含幺环，\((C,F)\) 为左 \(R\)-模的递减滤过上链复形。取 \(E_0^{p,q}=F^pC^{p+q}/F^{p+1}C^{p+q}\)，并对 \(r\ge1\) 以\cref{b4:eq:spectral-page-formula}定义 \(E_r^{p,q}\)。这些页在 \(\mathrm{d}_r[x]=[\mathrm{d}x]\) 下构成谱序列，且对滤过复形映射自然。相同结论适用于交换范畴中的递减滤过上链复形。
\end{theorem}
\begin{proof}
对 \(r\ge1\)，若 \(x\in Z_r^{p,q}\)，则 \(\mathrm{d}x\in F^{p+r}\) 且 \(\mathrm{d}^2x=0\)，所以它属于目标分子。若更改 \(x\) 一个 \(u\in Z_{r-1}^{p+1,q-1}\)，则 \(\mathrm{d}u\) 正好属于目标分母的第二项；更改一个余边不影响 \(\mathrm{d}x\)。故 \(\mathrm{d}_r\) 良定义，且平方为零。

现比较 \(H(E_r,\mathrm{d}_r)\) 与 \(E_{r+1}\)。若 \([x]\in\ker \mathrm{d}_r\)，可写
\[
 \mathrm{d}x=u+\mathrm{d}v,\qquad
 u\in Z_{r-1}^{p+r+1,q-r},\quad
 v\in Z_{r-1}^{p+1,q-1}.
\]
用 \(x-v\) 代替 \(x\) 不改变原类，而且 \(\mathrm{d}(x-v)=u\in F^{p+r+1}\)。所以每个新余圈都能由 \(Z_{r+1}^{p,q}\) 的元素表示。

再确定何时这种代表给出零。旧分母和新边界所产生的差具有形式
\[
 x=u+\mathrm{d}v+\mathrm{d}w,
\]
其中 \(u\in Z_{r-1}^{p+1,q-1}\)、\(v\in Z_{r-1}^{p-r+1,q+r-2}\)、\(w\in Z_r^{p-r,q+r-1}\)。若 \(x\in Z_{r+1}^{p,q}\)，则 \(\mathrm{d}u=\mathrm{d}x\in F^{p+r+1}\)，故 \(u\in Z_r^{p+1,q-1}\)。又 \(v+w\in F^{p-r}C^{n-1}\)，且 \(\mathrm{d}(v+w)=x-u\in F^pC^n\)，所以 \(v+w\in Z_r^{p-r,q+r-1}\)。因此核恰为
\[
 Z_r^{p+1,q-1}+\mathrm{d}Z_r^{p-r,q+r-1},
\]
正是第 \(r+1\) 页的分母。反方向包含可直接代入，得到所需同构。

滤过映射保持以上全部子对象，并与微分、商映射相容，所以诱导自然的页间映射。交换范畴中的交用拉回定义，像与商存在；上述每次表示差或选取原像均可沿满态射完成有限提升，故由广义元素正合性判据得到同样证明。
\end{proof}

构造的意义是逐步改进代表元：第 \(r\) 页的余圈不是已经满足 \(\mathrm{d}x=0\)，而是能够把 \(\mathrm{d}x\) 推到第 \(p+r\) 层。若 \(\mathrm{d}_r[x]\ne0\)，这种改进在此处遇到障碍；若它为零，则上一证明给出了下一次修正。

\ProseParagraphBreak
回到章首的四项双复形，\([a]\) 在第一页存留，而 \(\mathrm{d}a=c\) 仍落在中间列。用 \(-b\) 修正代表元后，\(\mathrm{d}(a-b)=-e\) 落入第二层，因而 \(\mathrm{d}_2[a]=-[e]\)；\secref{b4:sec:1405}会核对全部页与总复形。直接构造的原著表述可对照 Weibel\cite[Sections 5.2--5.4]{Weibel1994HomologicalAlgebra}；该书采用同调型指标，与本章固定的上同调型指标须作换算。
